\documentclass[12pt,a4paper]{article}
\usepackage{amsmath,amssymb,mathrsfs,tikz,times,pifont}
\usepackage[T1]{fontenc}
\usepackage{enumitem}
\usepackage{multicol}
\newcommand\circitem[1]{%
\tikz[baseline=(char.base)]{
\node[circle,draw=gray, fill=red!55,
minimum size=1.2em,inner sep=0] (char) {#1};}}
\newcommand\boxitem[1]{%
\tikz[baseline=(char.base)]{
\node[fill=cyan,
minimum size=1.2em,inner sep=0] (char) {#1};}}
\setlist[enumerate,1]{label=\protect\circitem{\arabic*}}
\setlist[enumerate,2]{label=\protect\boxitem{\alph*}}
%%%::::::by chnini ameur :::::::%%%
\everymath{\displaystyle}
\usepackage[left=1cm,right=1cm,top=1cm,bottom=1.7cm]{geometry}
\usepackage[colorlinks=true, linkcolor=blue, urlcolor=blue, citecolor=blue]{hyperref}
\usepackage{array,multirow}
\usepackage[most]{tcolorbox}
\usepackage{varwidth}
\usepackage{float} %pour utiliser l'option [H] qui force l'image à apparaître exactement à l'endroit où elle est placée dans le code.
\tcbuselibrary{skins,hooks}
\usetikzlibrary{patterns}
%%%::::::by chnini ameur :::::::%%%
\newtcolorbox{exa}[2][]{enhanced,breakable,before skip=2mm,after skip=5mm,
colback=yellow!20!white,colframe=black!20!blue,boxrule=0.5mm,
attach boxed title to top left ={xshift=0.6cm,yshift*=1mm-\tcboxedtitleheight},
fonttitle=\bfseries,
title={#2},#1,
% varwidth boxed title*=-3cm,
boxed title style={frame code={
\path[fill=tcbcolback!30!black]
([yshift=-1mm,xshift=-1mm]frame.north west)
arc[start angle=0,end angle=180,radius=1mm]
([yshift=-1mm,xshift=1mm]frame.north east)
arc[start angle=180,end angle=0,radius=1mm];
\path[left color=tcbcolback!60!black,right color = tcbcolback!60!black,
middle color = tcbcolback!80!black]
([xshift=-2mm]frame.north west) -- ([xshift=2mm]frame.north east)
[rounded corners=1mm]-- ([xshift=1mm,yshift=-1mm]frame.north east)
-- (frame.south east) -- (frame.south west)
-- ([xshift=-1mm,yshift=-1mm]frame.north west)
[sharp corners]-- cycle;
},interior engine=empty,
},interior style={top color=yellow!5}}
%%%%%%%%%%%%%%%%%%%%%%%

\usepackage{fancyhdr}
\usepackage{eso-pic}         % Pour ajouter des éléments en arrière-plan
% Commande pour ajouter du texte en arrière-plan
\usepackage{tkz-tab}
\AddToShipoutPicture{
    \AtTextCenter{%
        \makebox[0pt]{\rotatebox{80}{\textcolor[gray]{0.7}{\fontsize{5cm}{5cm}\selectfont PGB}}}
    }
}
\usepackage{lastpage}
\fancyhf{}
\pagestyle{fancy}
\renewcommand{\footrulewidth}{1pt}
\renewcommand{\headrulewidth}{0pt}
\renewcommand{\footruleskip}{10pt}
\fancyfoot[R]{
\color{blue}\ding{45}\ \textbf{2026}
}
\fancyfoot[L]{
\color{blue}\ding{45}\ \textbf{Prof:M. BA}
}
\cfoot{\bf
\thepage /
\pageref{LastPage}}
\begin{document}
\renewcommand{\arraystretch}{1.5}
\renewcommand{\arrayrulewidth}{1.2pt}
\begin{tikzpicture}[overlay,remember picture]
\node[draw=blue,line width=1.2pt,fill=purple,text=blue,inner sep=3mm,rounded corners,pattern=dots]at ([yshift=-2.5cm]current page.north) {\begingroup\setlength{\fboxsep}{0pt}\colorbox{white}{\begin{tabular}{|*1{>{\centering \arraybackslash}p{0.28\textwidth}} |*2{>{\centering \arraybackslash}p{0.2\textwidth}|} *1{>{\centering \arraybackslash}p{0.19\textwidth}|} }
\hline
\multicolumn{3}{|c|}{$\diamond$$\diamond$$\diamond$\ \textbf{Lycée de Dindéfélo}\ $\diamond$$\diamond$$\diamond$ }& \textbf{A.S. : 2025/2026} \\ \hline
\textbf{Matière: Mathématiques}& \textbf{Niveau : T}\textbf{S2} &\textbf{Date: 11/06/2026} & \textbf{Durée : 4 heures} \\ \hline
\multicolumn{4}{|c|}{\parbox[c]{10cm}{\begin{center}
\textbf{{\Large\sffamily Composition Du 2$ ^\text{\bf nd} $ Semestre}}
\end{center}}} \\ \hline
\end{tabular}}\endgroup};
\end{tikzpicture}

\vspace{3cm}
% ==================== EXERCICE 1 ====================
\section*{\underline{Exercice 1 : (5 pts)}}
\subsection*{\underline{Partie A}}
On considère le polynôme $P(x) = x^3 - 2x^2 - x + 2$.

\begin{enumerate}
    \item Calculer $P(1)$ et conclure. \hfill \textbf{0,25pt}
    
    \item Déterminer les réels a ; b et c tels que $P(x) = (x - 1)(ax^2 + bx + c)$. \hfill \textbf{0,5pt}
    
    \item On pose $Q(x) = x^2 - x - 2$
    \begin{enumerate}
        \item Donner une factorisation complète de $P(x)$ . \hfill \textbf{0,25pt}
        \item Résoudre dans $\mathbb{R}$  $P(x) = 0$ et $P(x) \leq 0$. \hfill \textbf{1pt}
        \item En déduire les solutions dans $\mathbb{R}$, l'équation $P(x-2) = 0$ et $P(x-2) \leq 0$. \hfill \textbf{1pt}
    \end{enumerate}
\end{enumerate}
\subsection*{\underline{Partie B}}
Soit l'équation $(E) : x^2 - 2x - 3 = 0$. Soit $x'$ et $x''$ les racines de $(E)$.

    Sans calculer $x'$ et $x''$,
    \begin{enumerate}
        \item Calculer : 
             \begin{multicols}{2}
    \setlength{\columnseprule}{0.1mm} % La largeur de la ligne verticale entre les colonnes
        \begin{enumerate}
        \item $A =x' + x''$\hfill \textbf{0,25pt}
        \item $B =x'x''$.\hfill \textbf{0,25pt}
        \end{enumerate}
  \end{multicols}
        \item En déduire les valeurs de 
     \begin{multicols}{3}
    \setlength{\columnseprule}{0.1mm} % La largeur de la ligne verticale entre les colonnes
        \begin{enumerate}
        \item $C = \dfrac{1}{x'} + \dfrac{1}{x''}$ \hfill \textbf{0,5pt}
        \item $D = {x'}^2 + {x''}^2$ \hfill \textbf{0,5pt}
        \item $E = {x'}^3 + {x''}^3$ \hfill \textbf{0,5pt}
        \end{enumerate}
     \end{multicols}
    \end{enumerate}

% ==================== EXERCICE 2 ====================
\section*{\underline{Exercice 2 : (3 pts)}}

Soit l'équation $(E) : x^2 - (2m + 1)x + m^2 - 2 = 0$

\begin{enumerate}
    \item Discuter suivant les valeurs de $m$, les solutions de l'équation $(E)$ \hfill \textbf{(0,75 pt)}
    
    \item Trouver une relation indépendante de $m$ entre $x'$ et $x''$ \hfill \textbf{(0,5 pt)}
    
    \item Déterminer $m$ pour que $(E)$ admette :
    \begin{enumerate}
        \item Deux solutions de signes contraires. \hfill \textbf{(0,25 pt)}
        \item Deux solutions positives. \hfill \textbf{(0,75 pt)}
        \item Deux solutions négatives. \hfill \textbf{(0,75 pt)}
    \end{enumerate}
\end{enumerate}
\newpage
% ==================== EXERCICE 3 ====================
\section*{\underline{Exercice 3 :(6 pts)} }
\begin{enumerate}
    \item Résoudre dans $\mathbb{R}$ les équations suivants :
    %\begin{multicols}{2}
    %\setlength{\columnseprule}{0.1mm} % La largeur de la ligne verticale entre les colonnes
    \begin{enumerate}
    \item $3x^2 - 5x + 2 = 0$\hfill \textbf{(0,25pt)}
    \item $2|x - 2|^2 + 7|x - 2| + 3 = 0$\hfill \textbf{(0,75 pt)}
    \item $\left(x + \dfrac{1}{x}\right)^2 = 4\left(x + \dfrac{1}{x}\right) - 3$\hfill \textbf{(0,75 pt)}
    \end{enumerate}
    %\end{multicols}
    \item Résoudre dans $\mathbb{R}$ inéquations suivants :
   \begin{enumerate}
     \item $x^2 - x + 6 \geq 0$\hfill \textbf{(0,25 pt)}
     \item $\dfrac{2x^2 + 3x - 5}{-x^2 + 3x + 10} \leq 0$\hfill \textbf{(0,75 pt)}
   \end{enumerate} 
        \item Résoudre dans $\mathbb{R}^2$ les systèmes suivants :
    %\begin{multicols}{2}
    %\setlength{\columnseprule}{0.1mm} % La largeur de la ligne verticale entre les colonnes 
    \begin{enumerate}
    \item $\begin{cases} x + y = 3 \\ x^2 + 3xy + y^2 = 11 \end{cases}$ \hfill \textbf{(0,75 pt)}
    \item $\begin{cases} x^2 + y^2 = 5 \\ xy = 2 \end{cases}$ \hfill \textbf{(0,75 pt)}
    \item $\begin{cases} x^3 + y^3 = 7 \\ x + y = 1 \end{cases}$ \hfill \textbf{(0,75 pt)}
    \end{enumerate}
    %\end{multicols}
    \item En utilisant la méthode de \textbf{Cramer}, résoudre le système : 
    $\begin{cases} 3x + 4y = 5 \\ 2x - 3y = 1 \end{cases}$ \hfill \textbf{(0,75 pt)}
\end{enumerate}

% ==================== EXERCICE 4 ====================
\section*{\underline{Exercice 4 :(6 pts)} }

Le plan est muni d'un repère orthonormé direct $(O; \vec{i}, \vec{j})$. Soit $(\mathcal{C})$ le cercle trigonométrique de centre $O$.

\begin{enumerate}
    \item Déterminer la mesure principale des angles orientés dont les mesures en radians sont données ci-dessous : \hfill \textbf{(2 pt)}
    \begin{enumerate}
        \item $x_1 = \frac{25\pi}{4}$
        \item $x_2 = -\frac{19\pi}{6}$
        \item $x_3 = \frac{2026\pi}{3}$
        \item $x_34 = -\frac{9\pi}{3}$
    \end{enumerate}

    \item Placer avec précision sur le cercle trigonométrique $(\mathcal{C})$ les points images $M_1$, $M_2$ et $M_3$ associés respectivement aux réels $x_1$, $x_2$ et $x_3$. \hfill \textbf{(2 pt)}

    \item Calculer les valeurs exactes du cosinus et du sinus de chacun de ces trois réels. \hfill \textbf{(2 pt)}

\end{enumerate}
\end{document}
